\section{Conventions}\label{sec:conventions}

Varieties are smooth, projective and connected over $\CC$ unless said otherwise. Cohomology is singular cohomology with rational coefficients, $H^{k}(X)=H^{k}(X,\QQ)$, with its Hodge structure. A Hodge class of degree $2p$ is an element of $H^{2p}(X)\cap H^{p,p}$; their space is $\Hdg^{p}(X)$. A class is algebraic if it is a rational combination of fundamental classes of subvarieties of codimension $p$; by Fulton \cite{Ful98}, intersection products, pull-backs and push-forwards along morphisms preserve algebraic classes. An algebraic correspondence from $Y$ to $X$ is an algebraic class $\Gamma$ on $Y\times X$, acting by $\Gamma_{*}\gamma=\mathrm{pr}_{X*}([\Gamma]\cup\mathrm{pr}_{Y}^{*}\gamma)$. The Hodge group $\Hg(V)$ of a rational Hodge structure $V$ is the smallest algebraic subgroup of $\mathrm{GL}(V)$ defined over $\QQ$ whose real points contain the circle $h(U(1))$; for a polarisable $V$ it is reductive, the Hodge classes in every tensor construction on $V$ are its invariants, and the Hodge group of a variety is that of its total cohomology \cite{Del82}. The Hodge conjecture for $X$ is the statement that every Hodge class on $X$ is algebraic; $\HC$ is the conjecture for all varieties and $\AB$ for all abelian varieties.

For an abelian variety $A$ of dimension $g$ with $H^{1}(A)=V$, the cohomology ring is $\bigwedge^{\bullet}V$. An abelian $2n$-fold $A$ is of Weil type for the imaginary quadratic field $K$ if $K$ acts on $A$ up to isogeny with $V$ a $K$-vector space of dimension $2n$ on which the two complex embeddings have equal multiplicity $n$ in $H^{1,0}$. The Weil classes are the elements of $\bigwedge^{2n}_{K}V\subseteq\bigwedge^{2n}V$, a two-dimensional rational space of Hodge classes of degree $2n$ \cite{Weil77}; $K^{\times}$ acts on it by $\tau\mapsto\tau^{2n}$, and $\eta^{k}$, for $\eta$ the class of a $K$-hermitian polarisation, carries $\tau\mapsto N_{K/\QQ}(\tau)^{k}$. For the general CM field we follow \cite{BMB26}: a family of Weil type for a CM field $F$ of degree $2m$ with $n\ge2$ and discriminant $\delta$ has a connected polarised period domain $\cD_{F}$ of dimension $n^{2}$, and its very general member has Hodge group $\SU(V,H)$ for the $F$-hermitian form $H$ of the polarisation.

A Mumford fourfold is a simple abelian fourfold $X$ whose Hodge group is the $\QQ$-form of $\mathrm{SL}_{2}^{3}$ given by the norm-one group $D^{1}$ of a quaternion algebra $D$ over a totally real cubic field $F$, split at one real place and definite at the other two \cite{Mum69}. Its Hodge ring is generated by the polarisation and one class of degree four that is not a polynomial in divisors, and $X\times X$ carries exceptional Hodge classes of degree four that are not products of divisor classes and are not Weil classes of any quotient \cite{MZ95,BMB26}. We write $T=D^{0}\subseteq D$ for the trace-zero part, a nine-dimensional $F$-line on which $\Hg(X)$ acts by conjugation; $T(-1)$ is a Hodge structure of K3 type with Hodge numbers $(1,7,1)$, and the exceptional classes of $X\times X$ are the real multiplication by $F$ on the copy of $T(-1)$ inside $H^{1}(X)\otimes H^{1}(X)$.

For a K3 surface $S$ the transcendental lattice $T(S)$ is the orthogonal complement of $\NS(S)$ in $H^{2}(S,\ZZ)$; $T(S)_{\QQ}$ is an irreducible Hodge structure, and $E=\End_{\Hdg}T(S)_{\QQ}$ is a totally real field or a CM field, acting by self-adjoint endomorphisms in the first case \cite{Zar83}. The K3 lattice is $\Lambda=U^{3}\oplus E_{8}(-1)^{2}$, and $\Lambda_{\QQ}=\Lambda\otimes\QQ$. For a rational quadratic space $W$ we use the discriminant $d(W)\in\QQ^{\times}/\QQ^{\times2}$ and the Hasse invariants $\varepsilon_{p}(W)=\prod_{i<j}(a_{i},a_{j})_{p}$ of a diagonalisation $\langle a_{1},\dots,a_{k}\rangle$, with $(\ ,\ )_{p}$ the Hilbert symbol; the classification of rational quadratic forms by rank, signature, discriminant and Hasse invariants, and the existence theorem for prescribed local invariants, are Serre's \cite[Chap.~IV]{Ser73}.

Perfect complexes on $A$ are objects of the bounded derived category of coherent sheaves; $\ch(\mathcal{E})\in H^{\mathrm{ev}}(A)$ is the Chern character, additive on triangles, multiplicative on exterior products, with $\ch(\mathcal{E}[1])=-\ch(\mathcal{E})$. The semiregularity map $\sigma_{\mathcal{E}}\colon\Ext^{2}(\mathcal{E},\mathcal{E})\to\bigoplus_{q}H^{q+2}(A,\Omega^{q})$ is that of Buchweitz and Flenner \cite{BF03}, after Bloch \cite{Blo72}; its injectivity implies that the Hodge locus of $\ch(\mathcal{E})$ coincides, near the given point, with the locus where $\mathcal{E}$ deforms \cite{BF03,Pri24}.

Throughout, an assertion attributed to a reference is cited and not reproved; everything else is proved in the text. The phrase ``very general'' means outside a countable union of proper closed analytic subsets of the parameter space.
